%---------- Inleiding ---------------------------------------------------------

% TODO: Is dit voorstel gebaseerd op een paper van Research Methods die je
% vorig jaar hebt ingediend? Heb je daarbij eventueel samengewerkt met een
% andere student?
% Zo ja, haal dan de tekst hieronder uit commentaar en pas aan.

%\paragraph{Opmerking}

% Dit voorstel is gebaseerd op het onderzoeksvoorstel dat werd geschreven in het
% kader van het vak Research Methods dat ik (vorig/dit) academiejaar heb
% uitgewerkt (met medesturent VOORNAAM NAAM als mede-auteur).
% 

\section{Inleiding}
\label{sec:inleiding}

Binnen de retailsector zijn niet alle producten even eenvoudig te identificeren. Voor veel verpakte producten kan gebruikgemaakt worden van een barcode of QR-code, maar bij verse producten zoals brood of groenten en fruit is dit niet altijd het geval. Hierdoor moet de klant of medewerker het product zelf correct herkennen en selecteren. Computer vision biedt mogelijkheden om dit proces gedeeltelijk te automatiseren door producten te herkennen op basis van camerabeelden.

\subsection{Kadering en casus}

Dit onderzoek vertrekt vanuit een concrete probleemsituatie bij \textbf{Delhaize Ledeberg}. Klanten nemen er zelf een brood uit het broodrek en kunnen dit vervolgens aan de broodmachine snijden en verpakken. Hiervoor zijn verschillende soorten broodzakken beschikbaar. Welke zak gebruikt moet worden, hangt af van het brood. Zo worden onder andere ronde en rechthoekige broden in verschillende zakken verpakt en wordt voor biobroden een afzonderlijke broodzak gebruikt.

In de praktijk kiezen klanten echter regelmatig een verkeerde broodzak. Hierdoor wordt het voor het kassapersoneel moeilijker om te bepalen welk brood zich in de zak bevindt, waardoor aan de kassa een verkeerd product kan worden geselecteerd. De behoefte van Delhaize Ledeberg bestaat daarom in een systeem dat de klant bij de broodsnijmachine kan helpen bij het kiezen van de juiste broodzak.

Een mogelijke oplossing is een computer-vision\-sys\-teem bestaande uit een camera, computer en scherm. Wanneer een klant een brood aan de broodsnijder plaatst, wordt een foto genomen waarop een AI-model het brood analyseert. Vervolgens kan op het scherm worden weergegeven welke broodzak de klant moet gebruiken.

\subsection{Probleemstelling}

Hoewel het bepalen van de juiste broodzak de directe praktische toepassing vormt, ontstaat hierbij een bredere technische uitdaging. Het assortiment van Delhaize Ledeberg bestaat uit ongeveer vijftig verschillende broodsoorten, waarvan verschillende broden visueel sterk op elkaar lijken. Het systeem moet daarom voldoende visuele kenmerken uit een foto kunnen halen om broden correct van elkaar te kunnen onderscheiden.

Vanuit deze probleemsituatie wordt de volgende centrale onderzoeksvraag opgesteld:

\begin{quote}
    \textit{In welke mate kunnen bestaande AI-modellen verschillende broodsoorten bij Delhaize Ledeberg accuraat herkennen op basis van foto's?}
\end{quote}

\subsection{Doelgroep}

De primaire doelgroep van dit onderzoek is \textbf{Delhaize Ledeberg}, met in het bijzonder de eindverantwoordelijken van de winkel. Voor deze bachelorproef is de co-promotor dan ook Denis Reyniers, Store Manager bij zowel Delhaize Ledeberg als Delhaize Sterre. Voor hem moet het onderzoek aantonen of automatische broodherkenning haalbaar is en in welke mate een dergelijk systeem ondersteuning kan bieden bij het kiezen van de juiste broodzak.

De resultaten kunnen daarnaast inzicht geven in een mogelijke verdere toepassing waarbij niet alleen de broodzakcategorie, maar ook de specifieke broodsoort automatisch wordt bepaald. Een effectieve integratie hiervan met bijvoorbeeld het kassasysteem of het afdrukken van een producttickets valt echter buiten de scope van deze bachelorproef.

\subsection{Onderzoeksdoelstelling}

Het doel van deze bachelorproef is om naast een theoretisch antwoord op de onderzoeksvraag ook een werkende \textit{Proof of Concept} (PoC) te ontwikkelen. Hiervoor wordt een dataset opgebouwd met foto's van broodsoorten uit het assortiment van Delhaize Ledeberg. Op basis daarvan worden bestaande AI-modellen getraind en geëvalueerd om na te gaan in welke mate zij de verschillende broden van elkaar kunnen onderscheiden.

Daarbij wordt onderzocht welke modellen geschikt zijn voor deze classificatietaak en welke invloed fine-tuning heeft op de nauwkeurigheid.

Het concrete eindresultaat is een PoC die op basis van een foto een brood analyseert en minstens kan bepalen tot welke broodzakcategorie het behoort. Daarnaast wordt onderzocht in welke mate hetzelfde systeem verder kan gaan en de specifieke broodsoort kan classificeren. Op basis van de behaalde resultaten wordt ten slotte een technisch onderbouwd advies geformuleerd over de haalbaarheid van automatische broodherkenning binnen de concrete context van Delhaize Ledeberg.

%---------- Stand van zaken ---------------------------------------------------

\section{Literatuurstudie}
\label{sec:literatuurstudie}

In dit hoofdstuk wordt een eerste overzicht gegeven van de literatuur die relevant is voor dit onderzoek. Aangezien dit een onderzoeksvoorstel betreft, ligt de focus op de bronnen die het probleem van automatische broodherkenning en mogelijke oplossingen onderbouwen. Daarbij wordt eerst gekeken naar bestaande toepassingen van computer vision op brood en bakkerijproducten. Vervolgens wordt ingegaan op de uitdaging om visueel sterk gelijkende producten van elkaar te onderscheiden en op technieken zoals transfer learning en fine-tuning. Tot slot wordt de link gelegd met automatische productherkenning binnen een retailomgeving.

\subsection{Het probleem: visuele herkenning van brood}

Computer vision wordt al gebruikt om verschillende visuele eigenschappen van brood automatisch te analyseren. \textcite{Olakanmi2023} geven een overzicht van bestaande beeldverwerkingssystemen voor bakkerijproducten en tonen aan dat kenmerken zoals kleur, vorm, grootte en textuur uit beelden kunnen worden afgeleid. Dergelijke kenmerken worden in hun onderzoek gebruikt om de kwaliteit en eigenschappen van brood op een objectieve manier te beoordelen. De auteurs wijzen er op dat de kwaliteit van de beeldopname beïnvloed kan worden door onder andere de gebruikte belichting en camera-op\-stel\-ling. Een consistente opnameomgeving is daarom belangrijk wanneer beelden automatisch geanalyseerd worden. Dit zijn zaken die we zeker mee kunnen nemen naar dit onderzoek.

Dat neurale netwerken subtiele verschillen tussen brood kunnen leren, werd reeds aangetoond in onderzoek naar de bruiningsgraad van brood. \textcite{Cotrim2020} gebruikten convolutionele neurale netwerken om verschillende stadia  verkleuring van broodkorsten automatisch te herkennen. Ook \textcite{Tantiphanwadi2022} onderzochten de classificatie van verschillende bruiningsstadia van brood en maakten hiervoor gebruik van een vooraf getraind VGG-16-model. Hoewel deze onderzoeken niet gericht zijn op het herkennen van verschillende broodsoorten, tonen ze wel aan dat deep-learning\-mo\-del\-len visuele verschillen tussen verwante broodbeelden kunnen leren.

Voor de casus van Delhaize Ledeberg ligt de uitdaging niet bij het bepalen van de kwaliteit of bruiningsgraad van één broodsoort maar bij het onderscheiden van vijftig verschillende broden. Sommige broden verschillen duidelijk in vorm, grootte of samenstelling, terwijl andere soorten visueel sterk op elkaar lijken. Hierdoor sluit het probleem aan bij \textit{fine-grained image classification}. Bij dergelijke classificatietaken moeten modellen onderscheid maken tussen subcategorieën die slechts kleine visuele verschillen vertonen, terwijl er tegelijkertijd variatie kan bestaan tussen verschillende voorbeelden van dezelfde klasse \autocite{Wei_2022}.

Deze uitdaging komt ook voor binnen automatische voedselherkenning. \textcite{Boyd_2024} onderzochten verschillende convolutionele neurale netwerken voor \textit{fine-grained} voedselclassificatie en tonen aan dat de keuze van het model een belangrijke invloed heeft op de uiteindelijke prestaties. Dit maakt het relevant om voor dit onderzoek niet op voorhand één model te kiezen maar verschillende bestaande modellen evalueren op dezelfde dataset.

\subsection{De oplossing: transfer learning en domeinspecifieke fine-tuning}

Het verzamelen van een gigantische gelabelde dataset voor ongeveer vijftig broodsoorten is binnen de scope van een Proof of Concept moeilijk haalbaar. Een mogelijke oplossing hiervoor is \textit{transfer learning}. Hierbij wordt kennis die een model tijdens een eerdere training heeft geleerd, hergebruikt voor een nieuwe maar gerelateerde taak. Volgens \textcite{Zhuang_2021} kan transfer learning de afhankelijkheid van grote hoeveelheden gelabelde data voor een nieuwe doeltaak verminderen.

Voor beeldclassificatie kan hiervoor vertrokken worden van modellen die reeds vooraf getraind zijn op grote beelddatasets. In plaats van een volledig model vanaf nul te trainen, kunnen de bestaande modelgewichten verder aangepast worden aan de nieuwe dataset. \textcite{Abou_Baker_2022} vergeleken verschillende neurale netwerken en verschillende strategieën voor transfer learning. Hun resultaten tonen aan dat de prestaties afhankelijk zijn van zowel de gekozen architectuur als de manier waarop het model verder wordt gefinetuned. Dit ondersteunt de keuze om voor deze bachelorproef meerdere bestaande modellen onder vergelijkbare omstandigheden te evalueren.

Ook specifiek voor brood bestaat hiervoor een relevant precedent. \textcite{Tantiphanwadi2022} gebruikten een vooraf getraind VGG-16-model voor de classificatie van broodbeelden met een relatief beperkte trainingsdataset. Zij combineerden transfer learning met fine-tuning en data-augmentatie om het model verder aan te passen aan hun broodbeelden. Hoewel hun classificatietaak beperkt bleef tot verschillende bruiningsstadia, toont dit aan dat transfer learning ook binnen het domein van broodbeelden toegepast kan worden wanneer slechts een beperkte hoeveelheid eigen trainingsdata beschikbaar is.

Voor de casus van Delhaize Ledeberg betekent dit dat een bestaand model als vertrekpunt kan dienen en vervolgens verder gefinetuned kan worden met foto's van de broodsoorten uit de winkel. Op die manier kan onderzocht worden in welke mate domeinspecifieke training noodzakelijk is om visueel sterk gelijkende broden correct van elkaar te onderscheiden.

\subsection{Productherkenning binnen een retailomgeving}

Naast specifiek onderzoek over brood is ook bestaande literatuur over automatische productherkenning in winkels relevant. \textcite{Ciapas_2023} onderzochten een systeem voor het automatisch classificeren van producten zonder barcode aan een zelfscankassa (Groenten, Fruit, Snoepjes, Donuts,...). Hun onderzoek toont aan dat in dergelijke omgevingen bijkomende problemen voorkomen, zoals producten die niet volledig zichtbaar zijn of beelden waarop het product niet duidelijk genoeg aanwezig is. Dit toont aan dat niet alleen het classificatiemodel, maar ook de omstandigheden waarin de beelden worden opgenomen een belangrijke rol spelen.

Binnen de Belgische retailsector bestaat bovendien een praktijkvoorbeeld van een gelijkaardige toepassing. Colruyt Group testte in 2019 een systeem waarbij een camera automatisch groenten en fruit aan de kassa herkende \autocite{ColruytGroup2019}. Het systeem stelde op basis van een foto een product voor, waarna dit door de klant of kassamedewerker bevestigd kon worden. Hoewel deze toepassing gericht was op groenten en fruit en niet op brood, toont ze aan dat cameragebaseerde productherkenning reeds binnen een Belgische supermarktcontext werd onderzocht.

De toepassing bij Delhaize Ledeberg verschilt echter van deze bestaande voorbeelden. Het primaire doel is niet om onmiddellijk een product aan de kassa af te rekenen, maar om de klant aan de broodsnijmachine te ondersteunen bij het kiezen van de juiste broodzak. Hiervoor hoeft het systeem in eerste instantie enkel voldoende informatie uit het broodbeeld te halen om de correcte zakcategorie te bepalen. De Proof of Concept gaat vervolgens een stap verder door te onderzoeken in welke mate ook de specifieke broodsoort automatisch geclassificeerd kan worden.

Uit de literatuur blijkt dus dat de afzonderlijke technische bouwstenen voor een dergelijk systeem al bestaan. Computer vision kan visuele eigenschappen van brood analyseren volgens \autocite{Olakanmi2023}, neurale netwerken kunnen subtiele verschillen tussen broodbeelden leren volgens \autocite{Cotrim2020,Tantiphanwadi2022}, en transfer learning kan gebruikt worden om bestaande modellen aan een nieuwe classificatietaak aan te passen volgens \autocite{Zhuang_2021,Abou_Baker_2022}. Daarnaast bestaan er al toepassingen van automatische productherkenning binnen een retailomgeving \autocite{Ciapas_2023,ColruytGroup2019}. Wat binnen de literatuur echter nog niet wordt behandeld, is de classificatie van een groot aantal visueel gelijkaardige commerciële broodsoorten binnen de concrete omgeving van een supermarkt. Deze combinatie vormt daarom het centrale onderzoeksgebied van deze bachelorproef.

%---------- Methodologie ------------------------------------------------------
\section{Methodologie}
\label{sec:methodologie}

De methodologie van dit onderzoek is opgebouwd rond het ontwikkelen en evalueren van een \textit{Proof of Concept} (PoC) voor automatische broodherkenning bij Delhaize Ledeberg. Het onderzoek wordt opgedeeld in drie grote fasen: het verzamelen van de nodige data, het trainen en evalueren van verschillende modellen en ten slotte het ontwikkelen en testen van de PoC.

\subsection{Fase 1: Dataverzameling en voorbereiding}
\emph{Tijdsindicatie: november 2026}

In de eerste fase wordt de literatuurstudie verder uitgebreid. Hierbij wordt gekeken welke modellen in gelijkaardige onderzoeken gebruikt worden en welke bestaande datasets van brood of gelijkaardige producten bruikbaar kunnen zijn voor transfer learning.

Daarnaast wordt een eigen dataset opgebouwd bij Delhaize Ledeberg. Het assortiment bestaat uit ongeveer vijftig verschillende broodsoorten. De literatuurstudie gaat bepalen hoeveel foto's er precies nodig zijn maar voor nu wordt er gestreefd naar 10 tot 20 verschillende fysieke exemplaren per broodsoort, waarvan telkens ongeveer vijf foto's worden genomen. DIt komt op een totaal aantal foto's van 2500 tot 5000. Hierbij wordt gevarieerd in onder andere de hoek, positie en belichting. De foto's worden tijdens het verzamelen gekoppeld aan de juiste broodsoort en broodzakcategorie.

De belangrijkste deliverables van deze fase zijn een gelabelde Delhaize-dataset, een selectie van bestaande datasets voor transfer learning en een selectie van AI-modellen die verder onderzocht zullen worden.

\subsection{Fase 2: Training en evaluatie}
\emph{Tijdsindicatie: eind november tot midden december 2026}

Eenmaal de dataset beschikbaar is, worden de geselecteerde modellen getraind en geëvalueerd. Hierbij wordt gebruikgemaakt van transfer learning en fine-tuning. Er wordt ook onderzocht of een extra trainingsstap op een bestaande dataset van brood of gelijkaardige producten een verbetering oplevert.

Omdat de dataset uit enkele duizenden foto's kan bestaan en meerdere modellen getraind moeten worden, is er hier voldoende rekenkracht voor nodig. Indien mijn persoonlijke laptop hiervoor niet volstaat, wordt onderzocht of trainingshardware beschikbaar is via het VIC (\textit{Virtual IT Company}) van HOGENT of via Azure of AWS.

De modellen worden onder dezelfde omstandigheden met elkaar vergeleken. Hierbij wordt enerzijds gekeken hoe accuraat ze de specifieke broodsoort kunnen classificeren en anderzijds of ze de juiste broodzakcategorie kunnen bepalen. Ook wordt bekeken welke broodsoorten vaak met elkaar verward worden.

De belangrijkste deliverables van deze fase zijn de getrainde modellen, de evaluatieresultaten en een selectie van het best presterende model voor de PoC.

\subsection{Fase 3: Proof of Concept en praktijktest}
\emph{Tijdsindicatie: midden december 2026 tot 8 januari 2027}

Het best presterende model wordt geïntegreerd in een eenvoudige PoC bestaande uit een camera, computer en scherm. Wanneer een brood voor de camera wordt geplaatst, wordt een foto genomen en toont het systeem welke broodzak gebruikt moet worden. Daarnaast wordt ook onderzocht of de specifieke broodsoort correct herkend kan worden.

Indien dit praktisch mogelijk is, wordt de PoC vervolgens getest bij de broodsnijmachine van Delhaize Ledeberg. Hierdoor kan nagegaan worden of de resultaten uit de testdataset ook standhouden in de echte winkelomgeving. In dezelfde periode worden de resultaten verwerkt en wordt de bachelorproef verder uitgeschreven.

De belangrijkste deliverables van deze fase zijn een werkende PoC, de resultaten van de praktijktest en een technisch onderbouwd advies voor Delhaize Ledeberg over de haalbaarheid van automatische broodherkenning.

%---------- Verwachte resultaten ----------------------------------------------
\section{Verwacht resultaat, conclusie}
\label{sec:verwachte_resultaten}

Het verwachte resultaat van dit onderzoek is een werkende Proof of Concept. De verwachting is dat vooral het bepalen van de juiste broodzakcategorie haalbaar zal zijn, aangezien de broden hiervoor slechts in een beperkt aantal categorieën moeten worden ingedeeld.

Het herkennen van de specifieke broodsoort zal waarschijnlijk moeilijker zijn aangezien er zoveel soorten zijn. Het onderzoek zal moeten uitwijzen met welke nauwkeurigheid deze van elkaar onderscheiden kunnen worden.

Voor Delhaize Ledeberg moet het onderzoek aantonen of automatische broodherkenning in de praktijk een oplossing kan bieden voor het foutief kiezen van broodzakken. Indien de PoC goede resultaten behaalt, kan dit een basis vormen voor een verdere uitwerking binnen Delhaize Ledeberg en eventueel andere winkels.

